\uuid{1iS3}
\titre{Étude de fonction}
\auteur{bodin ; hanani ; mzari}
\organisation{exo7}
\licence{CC BY-NC-SA 4.0}
\datecreate{2020-01-20}
\niveau{L1}
\module{Analyse}
\chapitre{Fonctions circulaires et hyperboliques inverses}
\sousChapitre{Autre}
\difficulte{3}
\theme{Fonctions usuelles, étude de fonction}
\exo7qcmid{lille-1.fonctions-usuelles.029}

\contenu{
\begin{question}
\qclassification{126.00}
\qtags{motcle=étude de fonction}
Soit $f(x)= x^x(1-x)^{1-x}$. On notera $D_f$ le domaine de définition de $f$.  Quelles sont les assertions vraies ?
\begin{answers}
    \good{$D_f=]0,1[$}

    \good{L'ensemble des valeurs de $f$ est $[\frac{1}{2},1[$.}

    \bad{$f$ est croissante $]0,1[$.}

    \good{$f$ est une bijection de $[\frac{1}{2},1[ $ dans $[\frac{1}{2},1[$.}
\end{answers}
\begin{explanations}
Par définition, $f(x)= \exp [x\ln x + (1-x)\ln (1-x)]$, on déduit que  $D_f=]0,1[$. $f$ est dérivable sur $]0,1[$ et $f'(x)= \ln (\frac{x}{1-x}) f(x)$. En étudiant les variations de $f$, $f$ n'est pas monotone sur  $]0,1[$. Elle établit une bijection de $]0,\frac{1}{2}]$ dans $[\frac{1}{2},1[$ et de $[\frac{1}{2},1[ $ dans $[\frac{1}{2},1[$.
\end{explanations}
\end{question}
}
